\section{TPM Device}\label{sec:Device Types / TPM Device}

The virtio TPM device provides a virtio device interface for accessing a
virtual TPM. It provides a mechanism to send a TPM command, which is to be
executed in a user-chosen TPM locality, and receive a reponse along with a
device status code.

\subsection{Device ID}\label{sec:Device Types / TPM Device / Device ID}
  46

\subsection{Virtqueues}\label{sec:Device Types / TPM Device / Virtqueues}
\begin{description}
\item[4]
\end{description}

\subsection{Feature bits}\label{sec:Device Types / TPM Device / Feature bits}

\subsection{Device configuration layout}\label{sec:Device Types / TPM Device / Device configuration layout}

\begin{lstlisting}
struct virtio_tpm_config {
        /* TPM's buffer size (r/o); typically ~4k */
        __virtio16 buffersize;
        /* TPM version; 1 = TPM 1.2, 2 = TPM 2 */
        __u8 tpm_version;
} __attribute__((packed));
\end{lstlisting}

\begin{description}
\item[\field{buffersize}] contains the TPM's buffersize. This is typically
around 4kb.

\item[\field{tpm_version}] contains the version of the TPM. '1' represents a
TPM 1.2 and '2' a TPM 2.
\end{description}

\subsection{Device Initialization}\label{sec:Device Types / TPM Device / Device Initialization}

\begin{enumerate}
\item The virtqueue is initialized
\item The driver SHOULD read the TPM version from the configuration space
\item The driver SHOULD read the TPM's buffer size from the configuration
space
\end{enumerate}

\subsection{Device Operation}\label{sec:Device Types / TPM Device / Device Operation}

\subsection{Device Operation: Request Queue}\label{sec:Device Types / TPM Device / Device Operation: Request Queue}

The driver enqueues a request consisting of the following 4 items:

\begin{enumerate}
\item A header
\item The buffer holding the TPM command
\item A sufficiently large buffer for the TPM response
\item A structure describing the success or failure of the operation.
\end{enumerate}

The following data structure describes the header:
\begin{lstlisting}
struct virtio_tpm_cmd_header {
	/* locality of the command */
	uint8_t locty;
} __attribute__((packed));
\end{lstlisting}

The \field{locty} allows a caller to specify the locality in which the TPM
command is supposed to be executed. Typically values from 0-3 are supported,
since locality 4 is reserved for hardware.

The following data structure holds the status code in the response:
\begin{lstlisting}
struct virtio_tpm_cmd_result {
	uint8_t status;
} __attribute__((packed));
\end{lstlisting}

The status code shows potential failure reasons for executing the TPM
command related to the buffers provided in the queue. The following status
codes are defined:

\begin{lstlisting}
#define VIRTIO_TPM_CMD_RESULT_SUCCESS      0
/* bad input such as buffers that are too short */
#define VIRTIO_TPM_CMD_RESULT_BAD_INPUT    1
\end{lstlisting}

\drivernormative{\subsubsection}{Driver Operation}{Device Types / TPM Device / Device Operation}

TPM device drivers are typically involved in interactions with the TPM for
power management. This involves sending a command to the TPM so that it
saves part of its current state to NVRAM (TPM 'Savestate'). On some
platforms, such as x86_64, it is the task of the firmware (BIOS/UEFI) to send a
command to the TPM for resuming (ACPI S3 resume) with the previously stored
state. To simplify the implementation across different architectures, a
virtio-tpm device driver MAY choose to send the necessary TPM Startup command
from the kernel device driver.

Further, a TPM device is typically initialized by the firmware of the
platform. The firmware may use the TPM device for implementing trusted boot,
which involves taking measurements (hashing) of firmware data and code,
extending platform configuration registers (PCRs) of the TPM, and
logging the measurements. The resulting log is typically made available to
OS kernel drivers through one of the following mechanisms:

\begin{enumerate}
\item The TPM2 ACPI table holds fields describing the localtion and size of the log
\item UEFI provides an API call for kernel drivers to get the log
\item Platforms providing an OpenFirmware device tree may provide an API to get the
log or make the log avaiable directly through the device tree (?).
\end{enumerate}

A platform kernel device driver implementor may choose any existing method
to get access to this log. This may include support for TPM 2 ACPI tables,
UEFI API call, etc.

\drivernormative{\subsubsection}{Device Operation}{Device Types / TPM Device / Device Operation}
